%% parts/p02/propositional-metatheory-supplement.tex --- retained editorial additions
\section*{Editorial supplement to the restored chapter}
\addcontentsline{toc}{section}{Editorial supplement to the restored chapter}
\noindent\textit{Reader note.} The following separately attributed editorial
additions are retained after the exact source chapter so that the inherited
order and bytes remain intact.  The Boolean Prime Ideal calibration belongs
with \cref{sec:prop-metatheory-03}, and the resolution-refutation DAG belongs
with \cref{sec:prop-metatheory-04}.  These additions do not replace exact
inherited statements, supply aliases for unlabeled source results, or change
any source or review status.

\begin{hbdefinition}{def:p02-arbitrary-prop-language}{Finitary propositional language on an arbitrary set}
For an arbitrary set \(P\), \(\operatorname{Form}(P)\) is the least set
containing \(P\) and closed under the fixed finitary connectives. A valuation
is a total function \(v:P\to\{0,1\}\), extended by structural recursion. The
word ``arbitrary'' is necessary because the source completeness proof
enumerates a countable alphabet. This stands as a set-indexed syntax/semantics,
differs only in cardinality rather than formula arity, and is used in the
choice-calibrated completeness theorem.  Write
\(\Gamma\vdash_{\mathsf{CPL}(P)}\varphi\) when there is a finite
discharge-labelled natural-deduction tree using exactly the inherited
\(\mathsf{ND}_{\mathsf{CPL}}\) rules, with formula metavariables ranging over
\(\operatorname{Form}(P)\) and every open assumption belonging to \(\Gamma\).
This is a \(P\)-parameterized proof relation, not a new rule system; it is used
in the theorem and its Lindenbaum quotient.
\index{propositional language!arbitrary alphabet}
\end{hbdefinition}

\begin{hbtheorem}{thm:p02-arbitrary-prop-compactness-bpi}{cited=HowardRubin1998}{Boolean Prime Ideal Theorem and unrestricted propositional compactness}
Working over \(\mathsf{ZF}\), the Boolean Prime Ideal Theorem (equivalently,
the Ultrafilter Lemma of Definition~\ref{def:ultrafilter-lemma}) is equivalent
to the following scheme: for every set \(P\) of propositional variables and
every set \(\Gamma\subseteq\operatorname{Form}(P)\), if each finite subset
of \(\Gamma\) is satisfiable, then \(\Gamma\) is satisfiable.
\end{hbtheorem}

See also \cite{Jech1973Choice,PaseauLeek2022} for the choice-principle and
logical formulations of this equivalence.  Theorem~\ref{thm:prop-compactness}
is weaker in set-theoretic strength: its fixed countable alphabet supports the
explicit enumeration used in Lemma~\ref{lem:prop-lindenbaum}, with no appeal
to a Boolean prime ideal.  The arbitrary-language theorem does not silently
upgrade that narrower source proof.

\begin{hbtheorem}{thm:p02-arbitrary-prop-completeness}{cited=HowardRubin1998}{Choice-calibrated propositional completeness}
Over \(\mathsf{ZF}\), the Boolean Prime Ideal Theorem proves: for every set
\(P\), set \(\Gamma\subseteq\operatorname{Form}(P)\), and formula \(\varphi\),
\[
 \Gamma\models\varphi\quad\Longrightarrow\quad
 \Gamma\vdash_{\mathsf{CPL}(P)}\varphi,
\]
where derivations remain finite and use only finitely many premises. In fact,
the corresponding unrestricted compactness/completeness scheme is equivalent
over \(\mathsf{ZF}\) to the Boolean Prime Ideal Theorem.
\end{hbtheorem}

\noindent\emph{Proof.}
Form the Lindenbaum Boolean algebra of formulas modulo provable equivalence and
let \(F\) be the proper filter generated by the classes of \(\Gamma\). If
\(\Gamma\nvdash_{\mathsf{CPL}(P)}\varphi\), then \([\varphi]\notin F\). The Boolean Prime Ideal
Theorem extends \(F\) to an ultrafilter avoiding \([\varphi]\). Its
characteristic homomorphism to the two-element Boolean algebra is a valuation
satisfying \(\Gamma\) and falsifying \(\varphi\), contradicting semantic
consequence. For a countable \(P\), the source's explicit enumeration replaces
this ultrafilter extension; no choice principle is silently imported into that
narrower proof.

\section{Propositional duality}
\label{sec:p02-propositional-duality}

\begin{hbdefinition}{def:p02-nnf-dual}{Propositional NNF duality}
Use the metalinguistic NNF presentation
\[
 \alpha ::= L\mid\bottom\mid\topv\mid(\alpha\land\beta)
             \mid(\alpha\lor\beta),
\]
where \(L\) is a literal.  The two units are displayed separately so their
swap is literal even though \(\topv\) is an abbreviation.  Define
\[
\begin{aligned}
 L^{\mathrm d}&=L,&
 \bottom^{\mathrm d}&=\topv,&
 \topv^{\mathrm d}&=\bottom,\\
 (\alpha\land\beta)^{\mathrm d}
   &=\alpha^{\mathrm d}\lor\beta^{\mathrm d},&&
 (\alpha\lor\beta)^{\mathrm d}
   =\alpha^{\mathrm d}\land\beta^{\mathrm d}.
\end{aligned}
\]
This is a metalinguistic operation on NNF presentations, not a new
object-language connective.  Its distinction from negating a formula is
needed to state the CNF--DNF and valid--unsatisfiable correspondences below.
\index{dual formula}\index{duality!propositional}
\end{hbdefinition}

\begin{hbtheorem}{thm:p02-propositional-duality}{proved=annex:p2.1.8}{Duality theorem}
Let \(\bar v(p)=1-v(p)\).  For every NNF presentation \(\alpha\),
\[
 v\entails\alpha^{\mathrm d}\quad\Longleftrightarrow\quad
 \bar v\not\entails\alpha.
\]
Consequently dualization is an involution, preserves logical equivalence,
exchanges valid formulas with unsatisfiable formulas, and exchanges CNF with
DNF.
\end{hbtheorem}

\begingroup\small
\begin{longtable}{@{}lll@{}}
\caption{The NNF duality dictionary.}
\label{tab:p02-duality}\\
\toprule Object & Dual & Consequence \\
\midrule\endfirsthead
\toprule Object & Dual & Consequence \\
\midrule\endhead
literal \(L\) & \(L\) & complemented valuation reverses its value \\
\(\bottom\) & \(\topv\) & false and true units swap \\
\(\land\) & \(\lor\) & outer conjunction becomes disjunction \\
\(\lor\) & \(\land\) & outer disjunction becomes conjunction \\
CNF & DNF & clauses become conjunctions of literals \\
valid & unsatisfiable & Theorem~\ref{thm:p02-propositional-duality} \\
\bottomrule
\end{longtable}
\endgroup

\section{Interpolation}
\label{sec:p02-prop-interpolation}

\begin{hbdefinition}{def:p02-interpolant}{Craig interpolant, propositional}
If \(\varphi\entails\psi\), a \emph{Craig interpolant} for that consequence
is a formula \(\theta\) such that
\(\varphi\entails\theta\), \(\theta\entails\psi\), and every variable of
\(\theta\) occurs in both \(\varphi\) and \(\psi\).  The name is needed to
distinguish this formula-level witness from an arbitrary intermediate
consequence: its common-vocabulary restriction is the nonredundant clause used
by the theorem and construction below.
\end{hbdefinition}

\begin{hbtheorem}{thm:p02-interpolation}{cited=Craig1957}{Propositional interpolation theorem}
Every semantic consequence \(\varphi\entails\psi\) in classical
propositional logic has a Craig interpolant.
\end{hbtheorem}

\begin{hbmethod}{def:p02-interpolant-truth-table}{Truth-table construction of a propositional interpolant}
Let \(X\) be the finite set of variables common to \(\varphi\) and \(\psi\),
ordered by variable index.
For each assignment \(u:X\to\{0,1\}\) extendible to a valuation satisfying
\(\varphi\), form the complete conjunction of \(X\)-literals describing
\(u\), and disjoin these conjunctions in lexicographic order with \(0<1\).  Empty
conjunction and disjunction use \(\topv\) and \(\bottom\), respectively.
The resulting formula uses only variables from \(X\); its two consequence
clauses are stated separately below so that this derived construction is not
attributed to Craig's original proof.
Part~7 gives the proof-theoretic Maehara construction.
\end{hbmethod}

\begin{hblemma}{thm:p02-interpolant-truth-table-correctness}{proved=annex:p2.3.9}{Correctness of the truth-table interpolant construction}
If \(\varphi\entails\psi\), the output of
Method~\ref{def:p02-interpolant-truth-table} is a Craig interpolant for that
consequence.
\end{hblemma}

\begin{figure}[tp]
\centering
\begin{tikzpicture}[
  clause/.style={draw,rounded corners,inner sep=4pt,font=\small},
  every edge/.style={draw,-{Latex[length=2mm]}},
  level distance=13mm,
  sibling distance=31mm]
\node[clause] (c1) at (0,0) {\(\{P,Q\}\)};
\node[clause] (c2) at (3,0) {\(\{\neg P,R\}\)};
\node[clause] (c4) at (6,0) {\(\{\neg R\}\)};
\node[clause] (c3) at (9,0) {\(\{\neg Q,R\}\)};
\node[clause] (np) at (4.5,-1.5) {\(\{\neg P\}\)};
\node[clause] (nq) at (7.5,-1.5) {\(\{\neg Q\}\)};
\node[clause] (q) at (2.25,-3) {\(\{Q\}\)};
\node[clause] (empty) at (4.875,-4.5) {\(\varnothing\)};
\draw[-{Latex[length=2mm]}] (c2) -- node[left,font=\scriptsize] {\(R\)} (np);
\draw[-{Latex[length=2mm]}] (c4) -- (np);
\draw[-{Latex[length=2mm]}] (c3) -- node[right,font=\scriptsize] {\(R\)} (nq);
\draw[-{Latex[length=2mm]}] (c4) -- (nq);
\draw[-{Latex[length=2mm]}] (c1) -- node[left,font=\scriptsize] {\(P\)} (q);
\draw[-{Latex[length=2mm]}] (np) -- (q);
\draw[-{Latex[length=2mm]}] (q) -- node[left,font=\scriptsize] {\(Q\)} (empty);
\draw[-{Latex[length=2mm]}] (nq) -- (empty);
\end{tikzpicture}
\caption{Resolution-refutation DAG.  Each two-parent merge is an instance of
the rule on \hyperref[sheet:resolution-cpl]{\texttt{resolution-cpl}}; the last node is the empty clause
.}
\label{fig:p02-resolution-dag}
\end{figure}

\section{Decision procedures and complexity}
\label{sec:p02-prop-decision-complexity}

\begin{hbtheorem}{thm:p02-decidability}{cited=EbbinghausFlumThomas1994}{Decidability of classical propositional logic}
There is a terminating algorithm which, on finite encoded inputs, decides
whether a propositional formula is satisfiable or valid, whether two formulas
are logically equivalent, and whether a finite set \(\Gamma\) semantically
entails a formula \(\varphi\).
\end{hbtheorem}

\begin{hbconvention}{def:p02-complexity-standing}{Complexity standing for propositional decision problems}
A decision problem is here a set \(L\subseteq\{0,1\}^*\).  The class
\(\mathsf P\) consists of the decision problems decided by a deterministic
multitape Turing machine in a number of steps bounded by a polynomial in the
input length.  A polynomial-time
many-one reduction from \(L\) to \(K\) is a polynomial-time computable map
\(f\) such that \(x\in L\) iff \(f(x)\in K\).  The class \(\mathsf{NP}\)
consists of those \(L\) for which some polynomial \(q\) and polynomial-time
predicate \(R\) satisfy, whenever \(n\) is the encoded length of \(x\),
\[
 x\in L\quad\text{iff}\quad
 \text{some string }y\text{ of length at most }q(n)\text{ satisfies }R(x,y).
\]
The class \(\mathsf{coNP}\) consists of complements of \(\mathsf{NP}\)
problems.  For a complexity class \(\mathcal C\), a problem \(K\) is
\(\mathcal C\)-complete under these reductions iff \(K\in\mathcal C\) and
every \(L\in\mathcal C\) has a polynomial-time many-one reduction to \(K\).
These metalinguistic definitions are stated locally because the complexity
claims use the distinctions now; Part~4 supplies the full class ladder.
\index{P@\(\mathsf P\)}\index{NP@\(\mathsf{NP}\)}
\index{coNP@\(\mathsf{coNP}\)}\index{polynomial-time many-one reduction}
\end{hbconvention}

\begin{hbdefinition}{def:p02-sat-taut}{SAT and TAUT}
Write each variable \(p_n\) as the character \texttt{p}, the ordinary binary
digits for \(n\), and the terminator \texttt{;}; use the fully parenthesized formula string;
and encode each character from the resulting finite alphabet by a fixed-width
binary block.  This metalinguistic encoding and its inverse are linear-time
and linearly intertranslatable with ordinary finite-alphabet formula encodings.
\(\mathsf{SAT}\) is the set of encodings of satisfiable formulas, and
\(\mathsf{TAUT}\) is the set of encodings of valid formulas.
Malformed bit strings belong to neither language; syntactic recognition is
linear-time for the fixed encoding.
These metalinguistic abbreviations are needed to distinguish existential
satisfiability from universal validity.  They are not the same decision
problem, although negation relates their complements, and both distinctions
are used in the next two results.
\index{SAT}\index{TAUT}
\end{hbdefinition}

\begin{hbtheorem}{thm:p02-sat-np-complete}{cited=AroraBarak2009}{SAT NP-completeness}
\(\mathsf{SAT}\) is \(\mathsf{NP}\)-complete under polynomial-time many-one
reductions.
\end{hbtheorem}

Cook's theorem-proving formulation is the historical source
\cite{Cook1971}; the status citation supports the modern formulation above.

\begin{hbcorollary}{thm:p02-taut-conp-complete}{cited=AroraBarak2009}{TAUT is coNP-complete}
\(\mathsf{TAUT}\) is \(\mathsf{coNP}\)-complete under polynomial-time
many-one reductions.
\end{hbcorollary}

The preceding classification does not settle whether
\(\mathsf P=\mathsf{NP}\) or \(\mathsf{NP}=\mathsf{coNP}\).  Part~4 locates
these open equalities in the complexity-class ladder; Part~12 treats SAT
solvers as algorithms rather than as complexity-class collapses.

\begin{hbdefinition}{def:p02-tractable-classes}{Horn, dual-Horn, 2-CNF, and XOR instances}
A clause is \emph{Horn} when it has at most one positive literal and
\emph{dual-Horn} when it has at most one negative literal.  A clause set is
\emph{2-CNF} when every clause has at most two literals.  An
\emph{XOR instance} is a finite system of linear equations over the
two-element field, presented as parity constraints on Boolean variables.
For a CNF input with no empty clause, let \(N\) denote its total number of
literal occurrences; an empty clause is detected first and makes the instance
unsatisfiable.  For a parity input, let \(m\) and \(n\) be its numbers of
equations and variables.
The corresponding satisfiability problems restrict their inputs to these
classes.  The four names distinguish nonidentical syntactic restrictions whose
different closure patterns license the algorithms stated below; the table uses
the distinction rather than treating them as synonyms for ``easy SAT.''
\index{Horn clause}\index{dual-Horn clause}\index{2-CNF}\index{XOR-SAT}
\end{hbdefinition}

\begin{hbtheorem}{thm:p02-horn-sat}{cited=DowlingGallier1984}{Horn-SAT tractability}
Horn satisfiability is decidable in time linear in the input size by forward
chaining (equivalently, unit propagation in the implicational presentation):
after the empty-clause check, the bound is \(O(N)\).
\end{hbtheorem}

\begin{hbcorollary}{thm:p02-dual-horn-sat}{proved=annex:p2.3.12}{Dual-Horn-SAT tractability}
Dual-Horn satisfiability is decidable in \(O(N)\) time.
\end{hbcorollary}

\begin{hbtheorem}{thm:p02-two-sat}{cited=AspvallPlassTarjan1979}{2-SAT tractability}
For a 2-CNF instance without an empty clause, form the implication graph on
literals by adding directed edges from \(\overline L\) to \(K\) and from
\(\overline K\) to \(L\) for every clause \(L\lor K\), treating a unit clause
\(L\) as \(L\lor L\).  The instance
is satisfiable iff no occurring variable \(p\) has both \(p\) and \(\neg p\)
in one strongly connected component.  The test runs in \(O(N)\) time.
\end{hbtheorem}

\begin{hbtheorem}{thm:p02-xor-sat}{cited=Schaefer1978}{XOR-SAT tractability}
XOR satisfiability is decidable by Gaussian elimination over the two-element
field, for example in \(O(mn\min\{m,n\})\) field operations with a dense
implementation.
\end{hbtheorem}

\begin{table}[tp]
\centering
\small
\begin{tabular}{@{}p{0.18\textwidth}p{0.37\textwidth}p{0.28\textwidth}@{}}
\toprule
Class & Recognition and decision method & Worst-case status \\
\midrule
Horn & At most one positive literal per clause; forward chaining & \(O(N)\) \\
Dual-Horn & At most one negative literal; complement variables and use Horn & \(O(N)\) \\
2-CNF & Clauses of size at most two; no literal shares an SCC with its complement & \(O(N)\) \\
XOR & Linear parity equations over the two-element field; Gaussian elimination & \(O(mn\min\{m,n\})\) dense field operations \\
\bottomrule
\end{tabular}
\caption{Tractable propositional satisfiability classes.}
\label{tab:p02-tractable-sat}
\end{table}

For detailed classical semantics and proof systems, see
\cite{EbbinghausFlumThomas1994,OpenLogic2026,Kleene1952}; for resolution and
normal-form transformations, see
\cite{Robinson1965,BaazEglyLeitsch2001}; and for complexity, see
\cite{Cook1971,AroraBarak2009}.  The first- and higher-order chapters give
their own topic-specific sources at the relevant definitions and theorems.
